% Generated draft: requires mathematical and source review.
\documentclass[11pt]{article}
\usepackage[margin=25mm]{geometry}
\usepackage{fontspec}
\setmainfont{Latin Modern Roman}
\usepackage{amsmath,amssymb,mathtools,unicode-math}
\setmathfont{Latin Modern Math}
\usepackage{enumitem,needspace,xurl,hyperref}
\setlength{\parindent}{0pt}
\setlength{\parskip}{5pt}
\setlength{\emergencystretch}{3em}
\newcommand{\N}{\mathbb{N}}
\newcommand{\Z}{\mathbb{Z}}
\newcommand{\Q}{\mathbb{Q}}
\newcommand{\R}{\mathbb{R}}
\newcommand{\C}{\mathbb{C}}
\newcommand{\F}{\mathbb{F}}
\newcommand{\A}{\mathbb{A}}
\newcommand{\PP}{\mathbb{P}}
\newcommand{\E}{\mathbb{E}}
\newcommand{\eps}{\varepsilon}
\newcommand{\dd}{\,\mathrm{d}}
\begin{document}
\section*{UnsolvedMath ID: 701468}
\textbf{Source problem code:} catalogue entry
\subsection{Definitions and mathematical statement}
Let \(G\) be a hyperbolic group, understood in the usual sense of a finitely generated group having a Gromov-hyperbolic Cayley graph. Let \(BG\) be a classifying space for \(G\), so that \(H_2(G;A)=H_2(BG;A)\) for \(A=\Z,\Q\). For a closed oriented surface \(S\) and a homomorphism \(\rho:\pi_1(S)\to G\), let \(f_\rho:S\to BG\) be a map inducing \(\rho\), and define
\[
c(S,\rho)=(f_\rho)_*[S]\in H_2(G;\Z).
\]
Write \(c_{\Q}(S,\rho)\) for its image under the coefficient map \(H_2(G;\Z)\to H_2(G;\Q)\).

\textbf{Conjecture (D. Calegari; supplied formulation).} Under the interpretation that the multiplier is a positive integer, for every hyperbolic group \(G\) and every \(\alpha\in H_2(G;\Q)\), there exist finitely many closed oriented surfaces \(S_i\), injective homomorphisms \(\rho_i:\pi_1(S_i)\to G\), and an integer \(n\geq 1\) such that
\[
\sum_i c_{\Q}(S_i,\rho_i)=n\alpha
\qquad\text{in }H_2(G;\Q).
\]
The surfaces, homomorphisms, and multiplier may depend on \(G\) and \(\alpha\).

\textbf{Unresolved source qualifications.} The supplied equation contains \(n\) without specifying its domain or explicitly requiring it to be nonzero. The positive-integer requirement above is an explicit interpretation requiring confirmation. The supplied statement imposes no genus restriction on the surfaces; whether the phrase ``closed surface subgroups'' is intended to require genus at least two remains unresolved. Here ``virtually represented'' is expressed by the displayed multiple-class equation, as in the supplied record; no passage to a finite-index subgroup of \(G\) is specified there.
\subsection{Short English statement}
The supplied conjecture asks whether every rational second homology class of a hyperbolic group has a positive integral multiple expressible as a finite sum of fundamental classes of closed oriented surfaces whose fundamental groups inject into the group. The multiplier convention and any intended genus restriction require confirmation.
\subsection{Sources}
\begin{itemize}
\item Local UnsolvedMath registry, supplied complete source record: permanent ID \(701468\), exact source problem code \texttt{catalogue entry}, title ``Kourovka Notebook Problem 18.39.'' The record attributes the problem to D. Calegari and gives the proposed year as 2014.
\item Supplied bibliographic attribution: \emph{The Kourovka Notebook (Unsolved Problems in Group Theory)}, 18th issue (2014), problem 18.39; also identified in the record with the 21st issue edition (2026), \texttt{arXiv:1401.0300v48}. The collection description names E. I. Khukhro and V. D. Mazurov as editors. This reference has not been independently consulted.
\item The local record labels the problem ``open'' and gives its update timestamp as \texttt{2026-10-16T14:00:00Z}. This is supplied registry metadata, not an independently verified mathematical or literature assessment. No separate literature-check date or literature-source links are supplied.
\end{itemize}
\end{document}
